% GENERATED FILE. Edit canonical lessons, then run npm run generate:books.
% canonical-source-hash: fnv1a64:cea691c15682ff0e
% canonical-lessons: HI-C215-jhutha, HI-C215-buddhiman, HI-C215-murkh, HI-C215-bahadur, HI-C215-dayalu

\chapter{False, Clever, Foolish, Brave, Kind}
\label{ch:hi-false-clever-foolish-brave-kind}
\hlchaptermodality{215}

\begin{chapteropening}
\textbf{By the end of this chapter:} \emph{I can say false, clever, foolish, brave and kind.}

Complete the last lesson of chapter 215: I can say false, clever, foolish, brave and kind.
\end{chapteropening}

\section[\texorpdfstring{\hi{झूठा} (jhūṭhā)}{झूठा (jhūṭhā)}]{\hi{झूठा} (jhūṭhā) — false, lying}

\label{lesson:HI-C215-jhutha}



\begin{quote}
Before the new one: say the Hindi for honest, then the Hindi for true.
\end{quote}

\subsection*{You'll want to know: \hi{झूठा}}
\textbf{\hi{झूठा}} — \emph{jhūṭhā} — \textquotedblleft{}false, lying\textquotedblright{}.

\begin{grammarlens}[title={What you've built}]
One more word for this chapter.
\end{grammarlens}

\subsection*{Your turn}
\begin{itemize}
\raggedright
  \item \emph{Say it:} \emph{jhūṭhā}
  \item \emph{Say it:} \emph{jhūṭhā}, once more
  \item \emph{Say it:} say \emph{jhūṭhā}
  \item \emph{Recall:} say the Hindi for honest, then the Hindi for true, then say \emph{jhūṭhā} again
\end{itemize}

\begin{tcolorbox}[breakable,colback=teal!4,colframe=teal!35!black,title={Before you move on}]
What does \hi{झूठा} mean? (\textquotedblleft{}False, lying\textquotedblright{}.) Say it once more.
\end{tcolorbox}
\section[\texorpdfstring{\hi{बुद्धिमान} (buddhimān)}{बुद्धिमान (buddhimān)}]{\hi{बुद्धिमान} (buddhimān) — clever, wise}

\label{lesson:HI-C215-buddhiman}



\begin{quote}
Before the new one: say the Hindi for true, then the Hindi for false.
\end{quote}

\subsection*{You'll want to know: \hi{बुद्धिमान}}
\textbf{\hi{बुद्धिमान}} — \emph{buddhimān} — \textquotedblleft{}clever, wise\textquotedblright{}.

\begin{grammarlens}[title={What you've built}]
One more word for this chapter.
\end{grammarlens}

\subsection*{Your turn}
\begin{itemize}
\raggedright
  \item \emph{Say it:} \emph{buddhimān}
  \item \emph{Say it:} \emph{buddhimān}, once more
  \item \emph{Say it:} say \emph{buddhimān}
  \item \emph{Recall:} say the Hindi for true, then the Hindi for false, then say \emph{buddhimān} again
\end{itemize}

\begin{tcolorbox}[breakable,colback=teal!4,colframe=teal!35!black,title={Before you move on}]
What does \hi{बुद्धिमान} mean? (\textquotedblleft{}Clever, wise\textquotedblright{}.) Say it once more.
\end{tcolorbox}
\section[\texorpdfstring{\hi{मूर्ख} (mūrkh)}{मूर्ख (mūrkh)}]{\hi{मूर्ख} (mūrkh) — foolish}

\label{lesson:HI-C215-murkh}



\begin{quote}
Before the new one: say the Hindi for false, then the Hindi for clever.
\end{quote}

\subsection*{You'll want to know: \hi{मूर्ख}}
\textbf{\hi{मूर्ख}} — \emph{mūrkh} — \textquotedblleft{}foolish\textquotedblright{}.

\begin{grammarlens}[title={What you've built}]
One more word for this chapter.
\end{grammarlens}

\subsection*{Your turn}
\begin{itemize}
\raggedright
  \item \emph{Say it:} \emph{mūrkh}
  \item \emph{Say it:} \emph{mūrkh}, once more
  \item \emph{Say it:} say \emph{mūrkh}
  \item \emph{Recall:} say the Hindi for false, then the Hindi for clever, then say \emph{mūrkh} again
\end{itemize}

\begin{tcolorbox}[breakable,colback=teal!4,colframe=teal!35!black,title={Before you move on}]
What does \hi{मूर्ख} mean? (\textquotedblleft{}Foolish\textquotedblright{}.) Say it once more.
\end{tcolorbox}
\section[\texorpdfstring{\hi{बहादुर} (bahādur)}{बहादुर (bahādur)}]{\hi{बहादुर} (bahādur) — brave}

\label{lesson:HI-C215-bahadur}



\begin{quote}
Before the new one: say the Hindi for clever, then the Hindi for foolish.
\end{quote}

\subsection*{You'll want to know: \hi{बहादुर}}
\textbf{\hi{बहादुर}} — \emph{bahādur} — \textquotedblleft{}brave\textquotedblright{}.

\begin{grammarlens}[title={What you've built}]
One more word for this chapter.
\end{grammarlens}

\subsection*{Your turn}
\begin{itemize}
\raggedright
  \item \emph{Say it:} \emph{bahādur}
  \item \emph{Say it:} \emph{bahādur}, once more
  \item \emph{Say it:} say \emph{bahādur}
  \item \emph{Recall:} say the Hindi for clever, then the Hindi for foolish, then say \emph{bahādur} again
\end{itemize}

\begin{tcolorbox}[breakable,colback=teal!4,colframe=teal!35!black,title={Before you move on}]
What does \hi{बहादुर} mean? (\textquotedblleft{}Brave\textquotedblright{}.) Say it once more.
\end{tcolorbox}
\section[\texorpdfstring{\hi{दयालु} (dayālu)}{दयालु (dayālu)}]{\hi{दयालु} (dayālu) — kind}

\label{lesson:HI-C215-dayalu}



\begin{quote}
Before the new one: say the Hindi for foolish, then the Hindi for brave.
\end{quote}

\subsection*{You'll want to know: \hi{दयालु}}
\textbf{\hi{दयालु}} — \emph{dayālu} — \textquotedblleft{}kind\textquotedblright{}.

\begin{grammarlens}[title={What you've built}]
One more word for this chapter.
\end{grammarlens}

\subsection*{Your turn}
\begin{itemize}
\raggedright
  \item \emph{Say it:} \emph{dayālu}
  \item \emph{Say it:} \emph{dayālu}, once more
  \item \emph{Say it:} say \emph{dayālu}
  \item \emph{Recall:} say the Hindi for foolish, then the Hindi for brave, then say \emph{dayālu} again
\end{itemize}

\begin{tcolorbox}[breakable,colback=teal!4,colframe=teal!35!black,title={Before you move on}]
What does \hi{दयालु} mean? (\textquotedblleft{}Kind\textquotedblright{}.) Say it once more.
\end{tcolorbox}
